\documentclass{beamer}
\setbeamerfont{headline}{size=\footnotesize}

\mode<presentation>{
\usetheme{Montpellier}
\setbeamercovered{transparent}
\usecolortheme{lsc}
}

\mode<handout>{
  % tema simples para ser impresso
  \usepackage[bar]{beamerthemetree}
   % Colocando um fundo cinza quando for gerar transparências para serem impressas
  % mais de uma transparência por página
  \beamertemplatesolidbackgroundcolor{black!5}
}


\usepackage{amsmath,amssymb}
\usepackage[brazil]{varioref}
\usepackage[english,brazil]{babel}
\usepackage[utf8]{inputenc}
\usepackage{fontawesome}
\usepackage{tikz}
%\usepackage[latin1]{inputenc}
\usepackage{graphicx}
\usepackage{listings}
\usepackage{url}
\usepackage{colortbl}
\usepackage{caption}
\usepackage{tikz}
\usepackage{transparent}
\usepackage{multirow}
\usepackage{graphicx}
\usepackage{listings}
\usepackage{color}
\usepackage{listings}
\usepackage{caption}
\usepackage{listings} %iclude code in your document

\usepackage{algpseudocode}

%% Custom Commands for Beamer Presentations
%% Este arquivo pode ser usado para definir comandos customizados

% Exemplo de comando customizado (descomente se necessário)
% \newcommand{\alert}[1]{\textcolor{red}{#1}}


\definecolor{maroon}{cmyk}{0, 0.87, 0.68, 0.32}
\definecolor{halfgray}{gray}{0.55}
\definecolor{ipython_frame}{RGB}{207, 207, 207}
\definecolor{ipython_bg}{RGB}{247, 247, 247}
\definecolor{ipython_red}{RGB}{186, 33, 33}
\definecolor{ipython_green}{RGB}{0, 128, 0}
\definecolor{ipython_cyan}{RGB}{64, 128, 128}
\definecolor{ipython_purple}{RGB}{170, 34, 255}

\usepackage{listings}
\lstset{
    breaklines=true,
    %
    extendedchars=true,
    literate=
    {á}{{\'a}}1 {é}{{\'e}}1 {í}{{\'i}}1 {ó}{{\'o}}1 {ú}{{\'u}}1
    {Á}{{\'A}}1 {É}{{\'E}}1 {Í}{{\'I}}1 {Ó}{{\'O}}1 {Ú}{{\'U}}1
    {à}{{\`a}}1 {è}{{\`e}}1 {ì}{{\`i}}1 {ò}{{\`o}}1 {ù}{{\`u}}1
    {À}{{\`A}}1 {È}{{\'E}}1 {Ì}{{\`I}}1 {Ò}{{\`O}}1 {Ù}{{\`U}}1
    {ä}{{\"a}}1 {ë}{{\"e}}1 {ï}{{\"i}}1 {ö}{{\"o}}1 {ü}{{\"u}}1
    {Ä}{{\"A}}1 {Ë}{{\"E}}1 {Ï}{{\"I}}1 {Ö}{{\"O}}1 {Ü}{{\"U}}1
    {â}{{\^a}}1 {ê}{{\^e}}1 {î}{{\^i}}1 {ô}{{\^o}}1 {û}{{\^u}}1
    {Â}{{\^A}}1 {Ê}{{\^E}}1 {Î}{{\^I}}1 {Ô}{{\^O}}1 {Û}{{\^U}}1
    {œ}{{\oe}}1 {Œ}{{\OE}}1 {æ}{{\ae}}1 {Æ}{{\AE}}1 {ß}{{\ss}}1
    {ç}{{\c c}}1 {Ç}{{\c C}}1 {ø}{{\o}}1 {å}{{\r a}}1 {Å}{{\r A}}1
    {€}{{\EUR}}1 {£}{{\pounds}}1
}

%%
%% Python definition (c) 1998 Michael Weber
%% Additional definitions (2013) Alexis Dimitriadis
%% modified by me (should not have empty lines)
%%
\lstdefinelanguage{Python}{
    morekeywords={access,and,break,class,continue,def,del,elif,else,except,exec,finally,for,from,global,if,import,in,is,lambda,not,or,pass,print,raise,return,try,while},%
    %
    % Built-ins
    morekeywords=[2]{abs,all,any,basestring,bin,bool,bytearray,callable,chr,classmethod,cmp,compile,complex,delattr,dict,dir,divmod,enumerate,eval,execfile,file,filter,float,format,frozenset,getattr,globals,hasattr,hash,help,hex,id,input,int,isinstance,issubclass,iter,len,list,locals,long,map,max,memoryview,min,next,object,oct,open,ord,pow,property,range,raw_input,reduce,reload,repr,reversed,round,set,setattr,slice,sorted,staticmethod,str,sum,super,tuple,type,unichr,unicode,vars,xrange,zip,apply,buffer,coerce,intern},%
    %
    sensitive=true,%
    morecomment=[l]\#,%
    morestring=[b]',%
    morestring=[b]",%
    %
    morestring=[s]{'''}{'''},% used for documentation text (mulitiline strings)
    morestring=[s]{"""}{"""},% added by Philipp Matthias Hahn
    %
    morestring=[s]{r'}{'},% `raw' strings
    morestring=[s]{r"}{"},%
    morestring=[s]{r'''}{'''},%
    morestring=[s]{r"""}{"""},%
    morestring=[s]{u'}{'},% unicode strings
    morestring=[s]{u"}{"},%
    morestring=[s]{u'''}{'''},%
    morestring=[s]{u"""}{"""}%
    %
    % {replace}{replacement}{lenght of replace}
    % *{-}{-}{1} will not replace in comments and so on
    literate=
    {á}{{\'a}}1 {é}{{\'e}}1 {í}{{\'i}}1 {ó}{{\'o}}1 {ú}{{\'u}}1
    {Á}{{\'A}}1 {É}{{\'E}}1 {Í}{{\'I}}1 {Ó}{{\'O}}1 {Ú}{{\'U}}1
    {à}{{\`a}}1 {è}{{\`e}}1 {ì}{{\`i}}1 {ò}{{\`o}}1 {ù}{{\`u}}1
    {À}{{\`A}}1 {È}{{\'E}}1 {Ì}{{\`I}}1 {Ò}{{\`O}}1 {Ù}{{\`U}}1
    {ä}{{\"a}}1 {ë}{{\"e}}1 {ï}{{\"i}}1 {ö}{{\"o}}1 {ü}{{\"u}}1
    {Ä}{{\"A}}1 {Ë}{{\"E}}1 {Ï}{{\"I}}1 {Ö}{{\"O}}1 {Ü}{{\"U}}1
    {â}{{\^a}}1 {ê}{{\^e}}1 {î}{{\^i}}1 {ô}{{\^o}}1 {û}{{\^u}}1
    {Â}{{\^A}}1 {Ê}{{\^E}}1 {Î}{{\^I}}1 {Ô}{{\^O}}1 {Û}{{\^U}}1
    {œ}{{\oe}}1 {Œ}{{\OE}}1 {æ}{{\ae}}1 {Æ}{{\AE}}1 {ß}{{\ss}}1
    {ç}{{\c c}}1 {Ç}{{\c C}}1 {ø}{{\o}}1 {å}{{\r a}}1 {Å}{{\r A}}1
    {€}{{\EUR}}1 {£}{{\pounds}}1
    %
    {^}{{{\color{ipython_purple}\^{}}}}1
    {=}{{{\color{ipython_purple}=}}}1
    %
    {+}{{{\color{ipython_purple}+}}}1
    *{-}{{{\color{ipython_purple}-}}}1
    {*}{{{\color{ipython_purple}$^\ast$}}}1
    {/}{{{\color{ipython_purple}/}}}1
    %
    {+=}{{{+=}}}1
    {-=}{{{-=}}}1
    {*=}{{{$^\ast$=}}}1
    {/=}{{{/=}}}1,
    %
    identifierstyle=\color{black}\ttfamily,
    commentstyle=\color{ipython_cyan}\ttfamily,
    stringstyle=\color{ipython_red}\ttfamily,
    keepspaces=true,
    showspaces=false,
    showstringspaces=false,
    %
    rulecolor=\color{ipython_frame},
    frame=single,
    frameround={t}{t}{t}{t},
    framexleftmargin=6mm,
    numbers=left,
    numberstyle=\tiny\color{halfgray},
    %
    %
    backgroundcolor=\color{ipython_bg},
    %   extendedchars=true,
    basicstyle=\scriptsize,
    keywordstyle=\color{ipython_green}\ttfamily,
}
% Adicione estas linhas ao seu preâmbulo em EstruturaGenerica.tex
\usepackage{tikz}

\usetikzlibrary{arrows.meta, shapes.geometric, positioning, calc, chains, shapes.multipart}

% Substitua o \tikzstyle antigo por este \tikzset, que é a sintaxe moderna
\tikzset{
    node_style/.style={
        rectangle split,
        rectangle split horizontal, % <--- AQUI ESTÁ A CORREÇÃO!
        rectangle split parts=3,
        draw=lscText,
        fill=lscgreen!70,
        rounded corners,
        text=lscText,
        font=\bfseries,
        minimum height=2em,
        text depth=0.5ex
    },
    pointer_style/.style={
        -{Stealth[length=3mm, width=2mm]},
        draw=lscgreenFonte,
        thick
    },
    label_style/.style={
        text=lscText,
        font=\small
    },
     % ADICIONE ESTE NOVO ESTILO AQUI
    stack_node_style/.style={
        draw=lscText,
        fill=lscgreen!70,
        rounded corners,
        text=lscText,
        font=\bfseries,
        minimum height=2em,
        minimum width=4em
    }
}
\beamertemplatetransparentcovereddynamic
\addtobeamertemplate{title page}{\centering{ \includegraphics[width=1.25cm]{figs/logo.png}}}{}
\title[Estrutura de Dados I]{Estruturas Genéricas}
\author[Elton Sarmanho]{ 
Professor: Elton Sarmanho\inst{1} 
\newline
E-mail: \href{mailto:eltonss@ufpa.br}{eltonss@ufpa.br}\newline
\href{https://github.com/eltonsarmanho}{\faGithub} \href{https://www.linkedin.com/in/elton-sarmanho-836553185/}{\faLinkedinSquare}
}   
  \institute[UFPA]{
   \inst{1}%
     Faculdade de Sistemas de Informação - UFPA/CUTINS
     }     
 
% Se comentar a linha abaixo, irá aparecer a data quando foi compilada a apresentação  
\date{\today}
\pgfdeclareimage[height= 1.25 cm]{inf}{figs/logo.png}
% pode-se colocar o LOGO assim
\logo{{\transparent{0.35}\pgfuseimage{inf}}}

% \AtBeginSection[]{
%   \begin{frame}<beamer>
%     \frametitle{Roteiro}
%     \tableofcontents[currentsection,currentsubsection]
%   \end{frame}  
% }
\setbeamertemplate{navigation symbols}{\insertframenavigationsymbol \insertsectionnavigationsymbol \insertbackfindforwardnavigationsymbol} 
\newcounter{saveenumi}
\newcommand{\seti}{\setcounter{saveenumi}{\value{enumi}}}
\newcommand{\conti}{\setcounter{enumi}{\value{saveenumi}}}

\expandafter\def\expandafter\insertshorttitle\expandafter{%
  \insertshorttitle\hfill%
  \insertframenumber\,/\,\inserttotalframenumber}
\begin{document}
\setbeamertemplate{logo}{}
\begin{frame}
\titlepage
\end{frame}
%Começar mostrar Logo apos slide de titulo
\logo{{\transparent{0.35}\pgfuseimage{inf}}}

\begin{frame}
\frametitle{Roteiro}
\tableofcontents[sections={1-2}]
\end{frame}
\begin{frame}
\frametitle{Roteiro}
\tableofcontents[sections={3-4}]
\end{frame}
\begin{frame}
\frametitle{Roteiro}
\tableofcontents[sections={5-6}]
\end{frame}

%Declaração das Imagens da Apresentação
\pgfdeclareimage[width= 3.5 in]{fig1}{figs/industriaJogos.png}%
\pgfdeclareimage[width= 3.5 in]{fig11}{figs/industriaJogos.png}%
\pgfdeclareimage[width= 3.5 in, height = 2.25 in]{fig2}{figs/arquiteturaGeral.png}%
%%%%%%%%%%%%%%%%%%%%%%%%%%%%%%%%%%%%%%%%
\section{Conceitos Gerais}

\subsection{Objetivos}
\begin{frame}{}
    
\begin{itemize}
 \item \hspace{-1 pt}Esta aula apresenta conceitos gerais pilhas, filas e listas. Ao final, você deverá compreender os seguintes tópicos:
 \begin{itemize}
  \item Representação visual dessas estruturas
  \item Representação computacional dessas estruturas
  \item Algoritmos de Pesquisa, Remoção e Inserção
  \item Aplicações Reais
 \end{itemize} 
\end{itemize}
\end{frame}

\section{Listas}
\subsection{Conceitos Gerais}

\begin{frame}
\begin{itemize}
 \item É uma estrutura de dados na qual os objetos estão organizados em ordem linear.
 \item Listas Estáticas
 \begin{itemize}
  \item representadas por vetores e índices.
 \end{itemize}
 \item Listas dinâmicas
 \begin{itemize}
  \item representadas por ponteiro (aponta para próximo elemento). Por isso o nome ``Lista encadeada''
  \item Toda lista encadeada tem pelo menos um ponteiro para o início.
 \end{itemize}

\end{itemize}


\end{frame}

\begin{frame}
\begin{itemize}
 \item A lista possui apenas um dado primitivo é dita \textbf{homogênea}
 \item A lista possui dados compostos é dita \textbf{heterogênea}
 \end{itemize}

\begin{center}
 \includegraphics[bb=0 0 777 353,scale=0.4]{figs/TipoLista.png}
 % TipoLista.png: 777x353 px, 72dpi, 27.41x12.45 cm, bb=0 0 777 353
\end{center}


\end{frame}

\begin{frame}
\begin{itemize}
 \item Podemos implementar as listas de forma estática e dinâmica. Existem 5 tipos:
\end{itemize}

\begin{description}
 \item[Lista simplesmente encadeada e não ordenada] 
 \item[Lista simplesmente encadeada e ordenada] 
 \item[Lista duplamente encadeada e não ordenada] 
 \item[Lista duplamente encadeada e ordenada] 
 \item[Lista circular] 
\end{description}


\end{frame}


\begin{frame}{Estrutura Geral - Lista dinâmica}
\begin{itemize}
    \item A ordem em uma lista encadeada é determinada por um ponteiro em cada objeto.
    \item Suporta diversas operações.
    \item \textbf{Lista duplamente encadeada}: Cada elemento possui uma chave e ponteiros para o nó anterior e o próximo.
\end{itemize}
\begin{center}
    \begin{tikzpicture}[node distance=1.5cm, start chain=going right]
        % 1. Nós da lista usando o novo estilo de 3 compartimentos
        %    O dado vai na 'nodepart{two}'
        \node[node_style, on chain] (n1) {\nodepart{two}9};
        \node[node_style, on chain] (n2) {\nodepart{two}16};
        \node[node_style, on chain] (n3) {\nodepart{two}4};

        % 2. Rótulo L.head
        \node[font=\bfseries, left=0.5cm of n1] (head_label) {L.head};

        % 3. Seta do L.head para o primeiro nó
        %    Aponta para o compartimento do meio (o dado)
        \draw[pointer_style] (head_label.east) -- (n1.two);

        % 4. Ponteiros "próximo" e "anterior"
        %    Conectam os compartimentos de ponteiro dos nós

        % Ponteiro n1 <-> n2
        % 'proximo' de n1 (n1.three) aponta para 'anterior' de n2 (n2.one)
        \draw[pointer_style] (n1.three) to[bend left=25] (n2.one);
        % 'anterior' de n2 (n2.one) aponta para 'proximo' de n1 (n1.three)
        \draw[pointer_style] (n2.one) to[bend left=25] (n1.three);

        % Ponteiro n2 <-> n3
        \draw[pointer_style] (n2.three) to[bend left=25] (n3.one);
        \draw[pointer_style] (n3.one) to[bend left=25] (n2.three);

    \end{tikzpicture}
\end{center}
\end{frame}



\begin{frame}{Estrutura Geral - Lista dinâmica}
\begin{itemize}
 \item Cada elemento da lista possui 3 atributos
 \begin{itemize}
  \item \textbf{chave}
  \item \textbf{ponteiro} para \textbf{próximo} elemento (sucessor)
  \item \textbf{ponteiro} para elemento \textbf{anterior} (predecessor)
 \end{itemize}
 
\end{itemize}
\begin{center}
    \begin{tikzpicture}[node distance=1.5cm, start chain=going right]
        % 1. Nós da lista usando o novo estilo de 3 compartimentos
        %    O dado vai na 'nodepart{two}'
        \node[node_style, on chain] (n1) {\nodepart{two}9};
        \node[node_style, on chain] (n2) {\nodepart{two}16};
        \node[node_style, on chain] (n3) {\nodepart{two}4};

        % 2. Rótulo L.head
        \node[font=\bfseries, left=0.5cm of n1] (head_label) {L.head};

        % 3. Seta do L.head para o primeiro nó
        %    Aponta para o compartimento do meio (o dado)
        \draw[pointer_style] (head_label.east) -- (n1.two);

        % 4. Ponteiros "próximo" e "anterior"
        %    Conectam os compartimentos de ponteiro dos nós

        % Ponteiro n1 <-> n2
        % 'proximo' de n1 (n1.three) aponta para 'anterior' de n2 (n2.one)
        \draw[pointer_style] (n1.three) to[bend left=25] (n2.one);
        % 'anterior' de n2 (n2.one) aponta para 'proximo' de n1 (n1.three)
        \draw[pointer_style] (n2.one) to[bend left=25] (n1.three);

        % Ponteiro n2 <-> n3
        \draw[pointer_style] (n2.three) to[bend left=25] (n3.one);
        \draw[pointer_style] (n3.one) to[bend left=25] (n2.three);

    \end{tikzpicture}
\end{center}

\end{frame}


\begin{frame}{Estrutura Geral - Lista dinâmica}
\begin{itemize}
 \item Dado $x$ um elemento da lista
 \begin{itemize}
  \item Se $x.anterior = null$ então é o primeiro elemento da lista
 \item Se $x.proximo = null$ então é o último elemento da lista
 \end{itemize}
\item $L.head$ ou $L.inicio$ que aponta para primeiro elemento da lista
\begin{itemize}
 \item $L.inicio = null$ a lista é vazia
\end{itemize}

\end{itemize}
\begin{center}
    \begin{tikzpicture}[node distance=1.5cm, start chain=going right]
        % 1. Nós da lista usando o novo estilo de 3 compartimentos
        %    O dado vai na 'nodepart{two}'
        \node[node_style, on chain] (n1) {\nodepart{two}9};
        \node[node_style, on chain] (n2) {\nodepart{two}16};
        \node[node_style, on chain] (n3) {\nodepart{two}4};

        % 2. Rótulo L.head
        \node[font=\bfseries, left=0.5cm of n1] (head_label) {L.head};

        % 3. Seta do L.head para o primeiro nó
        %    Aponta para o compartimento do meio (o dado)
        \draw[pointer_style] (head_label.east) -- (n1.two);

        % 4. Ponteiros "próximo" e "anterior"
        %    Conectam os compartimentos de ponteiro dos nós

        % Ponteiro n1 <-> n2
        % 'proximo' de n1 (n1.three) aponta para 'anterior' de n2 (n2.one)
        \draw[pointer_style] (n1.three) to[bend left=25] (n2.one);
        % 'anterior' de n2 (n2.one) aponta para 'proximo' de n1 (n1.three)
        \draw[pointer_style] (n2.one) to[bend left=25] (n1.three);

        % Ponteiro n2 <-> n3
        \draw[pointer_style] (n2.three) to[bend left=25] (n3.one);
        \draw[pointer_style] (n3.one) to[bend left=25] (n2.three);

    \end{tikzpicture}
\end{center}

\end{frame}




\begin{frame}{Estrutura Geral - Lista dinâmica}
\begin{itemize}
 \item Caso queira implementar uma lista simplesmente encadeada. Que devemos alterar ?
 \item Caso queira implementar uma lista circular. Que devemos alterar ?
 \item Caso queira implementar uma lista ordenada. Que devemos alterar ?
\end{itemize}
\begin{center}
    \begin{tikzpicture}[node distance=1.5cm, start chain=going right]
        % 1. Nós da lista usando o novo estilo de 3 compartimentos
        %    O dado vai na 'nodepart{two}'
        \node[node_style, on chain] (n1) {\nodepart{two}9};
        \node[node_style, on chain] (n2) {\nodepart{two}16};
        \node[node_style, on chain] (n3) {\nodepart{two}4};

        % 2. Rótulo L.head
        \node[font=\bfseries, left=0.5cm of n1] (head_label) {L.head};

        % 3. Seta do L.head para o primeiro nó
        %    Aponta para o compartimento do meio (o dado)
        \draw[pointer_style] (head_label.east) -- (n1.two);

        % 4. Ponteiros "próximo" e "anterior"
        %    Conectam os compartimentos de ponteiro dos nós

        % Ponteiro n1 <-> n2
        % 'proximo' de n1 (n1.three) aponta para 'anterior' de n2 (n2.one)
        \draw[pointer_style] (n1.three) to[bend left=25] (n2.one);
        % 'anterior' de n2 (n2.one) aponta para 'proximo' de n1 (n1.three)
        \draw[pointer_style] (n2.one) to[bend left=25] (n1.three);

        % Ponteiro n2 <-> n3
        \draw[pointer_style] (n2.three) to[bend left=25] (n3.one);
        \draw[pointer_style] (n3.one) to[bend left=25] (n2.three);

    \end{tikzpicture}
\end{center}

\end{frame}
\subsection{Operações}
\begin{frame}
\begin{itemize}
 \item Busca
  \item Inserção
  \item Remoção
  
\end{itemize}

\end{frame}



\begin{frame}[fragile]{Implementação em C}

\centering\footnotesize
\begin{lstlisting}[language=C]
#include <stdio.h>
#include <stdlib.h>

struct node {
   int chave;
   struct node *anterior;
   struct node *proximo;
};

struct node *head = NULL;//Nosso L.inicio ou L.head

\end{lstlisting}
\end{frame}


\begin{frame}[fragile]{Implementação em Python}

\centering\footnotesize
\begin{lstlisting}[language=Python]

class Nodo:
    def __init__(self, valor):
        self.chave = valor
        self.proximo = None #proximo
        self.anterior = None #anterior


class ListaDuplamenteEncadeada:
    def __init__(self):
        self.inicio = None;

\end{lstlisting}
\end{frame}


\begin{frame}{Operações: Busca}
\begin{itemize}
 \item O procedimento $List\_Search(L,K)$ realiza uma busca linear em busca da chave $k$ retornando um ponteiro para esse elemento
 \item Se nenhum objeto for encontrado, busca retorna $Null$
\end{itemize}


\end{frame}




\begin{frame}{Operações: Busca}
\begin{itemize}
 \item O procedimento $List\_Search(L,K)$ realiza uma busca linear em busca da chave $k$ retornando um ponteiro para esse elemento
 \item Se nenhum objeto for encontrado, busca retorna $Null$
\end{itemize}

\begin{center}
    \begin{tikzpicture}[node distance=1.5cm, start chain=going right]
        % 1. Nós da lista usando o novo estilo de 3 compartimentos
        %    O dado vai na 'nodepart{two}'
        \node[node_style, on chain] (n1) {\nodepart{two}9};
        \node[node_style, on chain] (n2) {\nodepart{two}16};
        \node[node_style, on chain] (n3) {\nodepart{two}4};

        % 2. Rótulo L.head
        \node[font=\bfseries, left=0.5cm of n1] (head_label) {L.head};

        % 3. Seta do L.head para o primeiro nó
        %    Aponta para o compartimento do meio (o dado)
        \draw[pointer_style] (head_label.east) -- (n1.two);

        % 4. Ponteiros "próximo" e "anterior"
        %    Conectam os compartimentos de ponteiro dos nós

        % Ponteiro n1 <-> n2
        % 'proximo' de n1 (n1.three) aponta para 'anterior' de n2 (n2.one)
        \draw[pointer_style] (n1.three) to[bend left=25] (n2.one);
        % 'anterior' de n2 (n2.one) aponta para 'proximo' de n1 (n1.three)
        \draw[pointer_style] (n2.one) to[bend left=25] (n1.three);

        % Ponteiro n2 <-> n3
        \draw[pointer_style] (n2.three) to[bend left=25] (n3.one);
        \draw[pointer_style] (n3.one) to[bend left=25] (n2.three);

    \end{tikzpicture}
\end{center}

\begin{itemize}
 \item $List\_Search(L,4)$ Retorna qual elemento ?
 \item $List\_Search(L,88)$ Retorna qual elemento ?
 \item Se existir dois elementos repetidos ? Qual retornar?
\end{itemize}

\end{frame}

\begin{frame}[fragile]
\frametitle{Algoritmo: Busca em Lista}
\begin{block}{Busca Linear (LIST-SEARCH)}
\begin{algorithmic}[1]
\Procedure{List-Search}{$L, k$}
    \State $x \gets L.inicio$ \Comment{Inicia a busca pelo primeiro nó}
    \While{$x \neq \text{NIL}$ \textbf{and} $x.chave \neq k$}
        \State $x \gets x.proximo$ \Comment{Avança para o próximo nó}
    \EndWhile
    \State \textbf{return} $x$ \Comment{Retorna o nó ou NIL se não encontrar}
\EndProcedure
\end{algorithmic}
\end{block}
\vfill
\begin{itemize}
    \item A busca percorre a lista sequencialmente a partir do início.
    \item \textbf{Complexidade (Pior Caso):} $O(n)$, pois pode ser necessário percorrer todos os $n$ elementos.
\end{itemize}
\end{frame}

\begin{frame}{Listas: Análise e Aplicações}
\footnotesize
\begin{block}{Análise de Complexidade (Lista Duplamente Encadeada com n elementos)}
    \centering
    \begin{tabular}{l c c c}
    \hline
    \textbf{Operação} & \textbf{Melhor Caso} & \textbf{Médio} & \textbf{Pior Caso} \\
    \hline
    Acesso / Busca   & $O(1)$ & $O(n)$ & $O(n)$ \\
    Inserção (início) & $O(1)$ & $O(1)$ & $O(1)$ \\
    Inserção (final)\* & $O(1)$ & $O(1)$ & $O(1)$ \\
    Remoção (início)  & $O(1)$ & $O(1)$ & $O(1)$ \\
    Remoção (meio)\*\* & $O(1)$ & $O(n)$ & $O(n)$ \\
    \hline
    \end{tabular}

\end{block}

\begin{columns}[T]
    \begin{column}{0.5\textwidth}
        \begin{alertblock}{Vantagens}
            \begin{itemize}
                \item Tamanho dinâmico, cresce conforme a necessidade.
                \item Inserção e remoção eficientes nas extremidades ($O(1)$).
            \end{itemize}
        \end{alertblock}
    \end{column}
    \begin{column}{0.5\textwidth}
        \begin{exampleblock}{Desvantagens}
            \begin{itemize}
                \item Acesso lento a elementos (requer travessia).
                \item Maior consumo de memória devido aos ponteiros.
            \end{itemize}
        \end{exampleblock}
    \end{column}
\end{columns}
\end{frame}


\begin{frame}[fragile]{Implementação em C}

\centering\footnotesize
\begin{lstlisting}[language=C]
int list_search(int k)
{
   struct node *x = head;//x = L.inicio
   int indice = 0;
   while(x != NULL) {
      if(x->chave == k)
        return indice;
      indice = indice + 1;
      x = x->proximo;
   }
   return -1;
}
\end{lstlisting}
\end{frame}

\begin{frame}[fragile]{Implementação em Python}

\centering\footnotesize
\begin{lstlisting}[language=Python]
    #Codigo Anterior

    def list_search(self,k):
        x = self.inicio;
        while x != None and x.chave != k:
            x = x.proximo;
        return x;

\end{lstlisting}
\end{frame}

\begin{frame}{Análise}
\begin{block}{Análise Assintótica}
 
 \begin{itemize}
  \item Considerando lista de tamanho $n$ 
  \begin{description}
  \item[Pior Caso] $O(n)$  
 \end{description}

 \end{itemize}

 \end{block}


\end{frame}


\begin{frame}{Operações: Inserção}
\begin{itemize}
 \item O procedimento $List\_Insert(L,x)$ insere um elemento \textbf{x} no início da lista
\end{itemize}

\begin{center}
 \includegraphics[bb=0 0 154 30,scale=2.2]{figs/PseudoCodeListaInsert.png}
 % PseudoCodeListaInsert.png: 641x127 px, 300dpi, 5.43x1.08 cm, bb=0 0 154 30
\end{center}

\end{frame}

\begin{frame}{Algoritmo: Inserção}
\begin{center}
 \includegraphics[bb=0 0 524 313,scale=0.5]{figs/PseudoCodeListaInsert2.png}
 % PseudoCodeListaInsert2.png: 524x313 px, 72dpi, 18.49x11.04 cm, bb=0 0 524 313
\end{center}


\end{frame}


\begin{frame}[fragile]{Implementação em C}

\centering\footnotesize
\begin{lstlisting}[language=C]
void list_insert(int value) {
   //Criacao do Nodo
   struct node *x = (struct node*) malloc(sizeof(struct node));
   x->chave = value;
   x->anterior = NULL;
   x->proximo = NULL;

   if(head == NULL) {
      head = x;//L.inicio = x
   } else {
      x->proximo = head;
      if(head != NULL)
        head->anterior = x;
      head = x;
      x->anterior = NULL;
   }
}

\end{lstlisting}
\end{frame}

\begin{frame}[fragile]{Implementação em Python}

\begin{lstlisting}[language=Python]
    #Codigo Anterior

    def list_insert(self,data):
        x = Nodo(data);
        # Se lista vazia entao inserir no Inicio
        if self.inicio is None:
            self.inicio = x
            return;

        x.proximo = self.inicio;
        if self.inicio != None:
            self.inicio.anterior = x;
        self.inicio = x
        x.anterior = None

\end{lstlisting}
\end{frame}


\begin{frame}[fragile]{Implementação em C}

\centering\footnotesize
\begin{lstlisting}[language=C]
void print_list() {
   struct node *temp = head;
   while(temp != NULL) {
      printf("%d ", temp->chave);
      temp = temp->proximo;
   }
}

\end{lstlisting}
\end{frame}


\begin{frame}[fragile]{Implementação em Python}

\centering\footnotesize
\begin{lstlisting}[language=Python]

    def print_list(self):
        current = self.inicio
        while current:
            print(current.chave)
            current = current.proximo



\end{lstlisting}
\end{frame}



\begin{frame}{Análise}
\begin{block}{Análise Assintótica}
 
 \begin{itemize}
  \item O tempo de execução para $List\_Insert(L,x)$ inserir um elemento na lista de tamanho $n$.
  \begin{description}
  \item[Pior Caso] $O(1)$  
 \end{description}
\item Qual o tempo de execução para uma lista de tamanho $n$, caso a inserção fosse no final da lista ?
 \end{itemize}

 \end{block}


\end{frame}

\begin{frame}{Operações: Remoção}
\begin{itemize}
 \item O procedimento $List\_Delete(L)$ remove um elemento \textbf{x} no início da lista
\end{itemize}
\end{frame}


\begin{frame}{Operações: Remoção}
\begin{itemize}
 \item O procedimento $List\_Delete(L)$ remove um elemento \textbf{x} no início da lista
\end{itemize}
\begin{center}
 \includegraphics[bb=0 0 171 109,scale=1.4]{figs/RemoveElementList.png}
 % RemoveElementList.png: 712x454 px, 300dpi, 6.03x3.84 cm, bb=0 0 171 109
\end{center}


\end{frame}



\begin{frame}{Análise}
\begin{block}{Análise Assintótica}
 
 \begin{itemize}
  \item O tempo de execução para $List\_Delete(L,x)$ remover um elemento na lista de tamanho $n$.
  \begin{description}
  \item[Melhor Caso] $O(1)$
  \item[Pior Caso] $O(n)$
 \end{description}
 \end{itemize}

 \end{block}


\end{frame}




\begin{frame}{Algoritmo: Remoção}
\begin{center}
 \includegraphics[bb=0 0 66 34,scale=3.2]{figs/PseudoCodeListaRemove.png}
 % PseudoCodeListaRemove.png: 275x141 px, 300dpi, 2.33x1.19 cm, bb=0 0 66 34
\end{center}
\end{frame}

\begin{frame}{Algoritmo: Remoção}
\begin{center}
 \includegraphics[bb=0 0 911 462,scale=0.3]{figs/PseudoCodeListaRemove_2.png}
 % PseudoCodeListaRemove_2.png: 911x462 px, 72dpi, 32.14x16.30 cm, bb=0 0 911 462
\end{center}

\end{frame}


\begin{frame}[fragile]{Implementação em C}

\centering\footnotesize
\begin{lstlisting}[language=C]
void list_delete(int x)
{
   struct node *current = head;//x = L.inicio
   while(current != NULL) {
      if(current->chave == x)
      {
        if(current->anterior != NULL)
            current->anterior->proximo = current->proximo;
        else head = current->proximo;
        if (current->proximo != NULL)
            current->proximo->anterior = current->anterior;
      }//Fim do IF
    current = current->proximo;
   }//Fim do While
}//Fim do Metodo
\end{lstlisting}
\end{frame}

\begin{frame}[fragile]{Implementação em Python}
\centering\footnotesize
\begin{lstlisting}[language=Python]
    #Codigo Anterior

def list_delete(self,x):
  current = self.inicio;
  while current != None:
    if current.chave == x:
      if current.anterior is not None:
        current.anterior.proximo = current.proximo
      else:
        self.inicio = current.proximo
      if current.proximo is not None:
        current.proximo.anterior = current.anterior
      return
  current = current.proximo;

\end{lstlisting}
\end{frame}


\section{Estrutura de Dados: Pilha e Fila}
\subsection{Características Gerais}
\begin{frame}
 \begin{itemize}
  \item São consideradas listas especializadas por possuírem características próprias.
  \item Ambas possuem operações de próprias de inserir, excluir, encontrar o maior elemento e contar os elementos.
  \item Dados organizados em ordem linear
  \item Quando representadas por vetores são consideradas \textbf{estáticas}. Quando representadas por listas encadeadas com ponteiros, são consideradas \textbf{dinâmicas}
  \item Quando formado por dado primitivo, dizemos que são estruturas \textbf{homogênea}. Quando representadas por dados compostos, dizemos que são \textbf{heterogênea}.
 \end{itemize}

\end{frame}



\section{Pilha (Stack)}
\subsection{Conceitos Gerais}

\begin{frame}{Pilha (Stack): Conceitos Gerais}
\begin{itemize}
    \item Uma coleção ordenada de itens onde as operações ocorrem na mesma extremidade: o \textbf{topo}.
    \item Implementa a política \textbf{LIFO (Last-In, First-Out)}: o último a entrar é o primeiro a sair.
\end{itemize}
\begin{center}
\begin{tikzpicture}
    % Representação da Pilha (cilindro)
    \draw[thick, draw=lscText] (-1,0) -- (-1,3) -- (1,3) -- (1,0);

    % Elementos da pilha usando o novo estilo 'stack_node_style'
    \node[stack_node_style] at (0,0.5) {1};
    \node[stack_node_style] at (0,1.5) {2};

    % reticências para indicar mais elementos
    \node[label_style, fill=white] at (0,2.5) {...};

    % Setas e Rótulos
    \draw[pointer_style, red] (-1.2, 2.5) -- node[left] {Push (Empilhar) \hspace*{0.5cm}} (-2, 2.5);
    \draw[pointer_style, red] (1.2, 2.5) -- node[right] {\hspace*{0.5cm}Pop (Desempilhar)} (2, 2.5);
    \node[label_style, font=\bfseries, red] at (0,3.3) {Topo};
\end{tikzpicture}
\end{center}
\end{frame}

\begin{frame}{Comportamento}

\begin{center}
 \includegraphics[bb=0 0 120 73,width = \linewidth]{figs/ComportamentoPilha.png}
 % ComportamentoPilha.png: 498x305 px, 300dpi, 4.22x2.58 cm, bb=0 0 120 73
\end{center}

\end{frame}


\begin{frame}{Análise}
 \begin{block}{Análise da Complexidade}
 \begin{description}
  \item[Push/Pop] Operação de tempo constante, então temos $O(1)$
  \item[Consulta] No pior caso, temos $O(n)$. No melhor caso temos $O(1)$
 \end{description}

 \end{block}


\end{frame}

\begin{frame}{Algoritmo}
\begin{center}
 \includegraphics[bb=0 0 360 440,width = 0.45\linewidth]{figs/PseudoCodePilha.png}
 % PseudoCodePilha.png: 360x477 px, 72dpi, 12.70x16.83 cm, bb=0 0 360 477
\end{center}

\end{frame}


\begin{frame}[fragile]{Implementação em C}

\centering\footnotesize
\begin{lstlisting}[language=C]
#include <stdio.h>
#include <stdlib.h>
#include<stdbool.h>
#define MAX_SIZE 10

int stack[MAX_SIZE];
int top = -1;

bool stack_empty()
{
    if(top == 0)
        return true;
    else return false;
}
\end{lstlisting}
\end{frame}


\begin{frame}[fragile]{Implementação em C}

\centering\footnotesize
\begin{lstlisting}[language=C]
void push(int x) {
    if (top == MAX_SIZE - 1) {printf("Pilha cheia\n");
        return;
    }
    top++;
    stack[top] = x;
}
int pop() {
    if (stack_empty()) {printf("Pilha vazia\n");
        return -1;
    }
    else
    {
        top--;
        return stack[top+1];
    }
}


\end{lstlisting}
\end{frame}

\begin{frame}[fragile]{Implementação em Python}

\centering\footnotesize
\begin{lstlisting}[language=Python]
class Stack:
    def __init__(self):
        self.stack = []
        self.top = -1

    def stack_empty(self):
        if self.top == -1:
            return True;
        else:
            return False;

\end{lstlisting}
\end{frame}


\begin{frame}[fragile]{Implementação em Python}

\centering\footnotesize
\begin{lstlisting}[language=Python]
class Stack:

    # Codigo Anterior
    def push(self, item):
        self.top += 1
        self.stack.insert(self.top,item)

    def pop(self):
        if self.stack_empty():
            return None
        item = self.stack[self.top]
        self.top -= 1
        self.stack.pop()
        return item

\end{lstlisting}
\end{frame}



\begin{frame}{Aplicações}

\begin{itemize}
 \item \textit{Undo/Redo} no editor de textos
 \item Controle de chamada de funções
\end{itemize}
\begin{center}
 \includegraphics[bb=0 0 168 39,scale=1.7]{figs/ChamadaFuncoes.png}
 % ChamadaFuncoes.png: 700x161 px, 300dpi, 5.93x1.36 cm, bb=0 0 168 39
\end{center}

\end{frame}





\section{Fila (Queue)}



\subsection{Conceitos Gerais}

\begin{frame}
\footnotesize
 \begin{itemize}
 \item Uma fila é uma coleção ordenada de zero ou mais itens, de um mesmo tipo ou não, tal que suas operações principais, \underline{inserção} e \underline{remoção}, são realizadas nas extremidades, denominadas \textbf{INICIO (CABEÇA)} e \textbf{FIM (CAUDA)} .
  \item Implementa a política de\textbf{ primeiro a entrar, primeiro a sair} ou FIFO (\textit{first-in, first-out}).
 \end{itemize}
\begin{center}
    % O único ajuste é trocar 'node_style' por 'stack_node_style'
    \begin{tikzpicture}[node distance=0.3cm]
        \node[label_style] (enqueue_label) {Enqueue};
        % Nó que vai entrar na fila
        \node[stack_node_style, left=0cm of enqueue_label, fill=lscgreenFonte!70] (new) {};

        % Nós que representam a fila

        \node[stack_node_style,right=0.2cm of enqueue_label,yshift=-1cm] (back) {};
        \node[stack_node_style, right=of back] (n1) {};
        \node[stack_node_style, right=of n1] (front) {};

        % Nó que vai sair da fila
        \node[stack_node_style, right= 1cm of front, fill=lscgreenFonte!70,yshift=1cm] (old) {};
        \node[label_style, right=0.1cm of old] {Dequeue};

        \draw[pointer_style, red] (new) -- (back);
        \draw[pointer_style, red] (front) -- (old);

        \node[label_style, above=0.3cm of back] {Final};
        \node[label_style, above=0.3cm of front] {Início};
    \end{tikzpicture}
\end{center}


\end{frame}

\begin{frame}
 \begin{block}{Operações}
  \begin{description}
   \item[\textit{Enqueue} (Enfileirar)] Insere um novo item no \textbf{início} da estrutura
   \item[\textit{Dequeue} (Desinfileirar)] Remove um item no \textbf{fim} da estrutura
  \end{description}
 \end{block}



\end{frame}


\begin{frame}{Comportamento}

\begin{description}
 \item[Início] aponta para o primeiro da fila, ou seja, o
primeiro elemento a sair
\item [Fim] aponta para o fim da fila, ou seja, onde o
próximo elemento entrará
\end{description}

\end{frame}

\begin{frame}{Comportamento}

\begin{center}
 \includegraphics[bb=0 0 112 55,width=0.6\textwidth]{figs/ComportamentoFila_2.png}
 % ComportamentoFila_2.png: 468x228 px, 300dpi, 3.96x1.93 cm, bb=0 0 112 55
\end{center}


\end{frame}

\begin{frame}{Comportamento}
\begin{itemize}
 \item $Enqueue(A)\rightarrow Enqueue(B)\rightarrow Enqueue(C)$
\end{itemize}


\begin{center}
 \includegraphics[bb=0 0 109 55,width=0.6\textwidth]{figs/ComportamentoFila_3.png}
 % ComportamentoFila_3.png: 456x231 px, 300dpi, 3.86x1.96 cm, bb=0 0 109 55
\end{center}


\end{frame}

\begin{frame}{Comportamento}
\begin{itemize}
 \item $Enqueue(Z)\rightarrow Enqueue(R)\rightarrow Enqueue(S)$
\end{itemize}


\begin{center}
 \includegraphics[bb=0 0 112 55,width=0.6\textwidth]{figs/ComportamentoFila_4.png}
 % ComportamentoFila_4.png: 467x228 px, 300dpi, 3.95x1.93 cm, bb=0 0 112 55
\end{center}



\end{frame}


\begin{frame}{Comportamento}
\begin{itemize}
 \item $Dequeue()\rightarrow Dequeue()$
\end{itemize}

\begin{center}
 \includegraphics[bb=0 0 112 55,width=0.6\textwidth]{figs/ComportamentoFila_5.png}
 % ComportamentoFila_4.png: 467x228 px, 300dpi, 3.95x1.93 cm, bb=0 0 112 55
\end{center}



\end{frame}


\begin{frame}{Comportamento}
\begin{itemize}
 \item Temos algumas questões ?
 \begin{itemize}
  \item Como reutilizar os espaços do início da fila?
  \item Como inserir mais elementos?
 \end{itemize}

\end{itemize}

\begin{center}
 \includegraphics[bb=0 0 112 55,width=0.6\textwidth]{figs/ComportamentoFila_5.png}
 % ComportamentoFila_4.png: 467x228 px, 300dpi, 3.95x1.93 cm, bb=0 0 112 55
\end{center}

\end{frame}

\begin{frame}{Comportamento: Fila Circular}
\footnotesize
\begin{itemize}
  \item A \textbf{Fila Circular} otimiza o uso de um vetor, tratando-o como um círculo.
  \item Permite reutilizar os espaços liberados no início do vetor.
  \item O cálculo dos índices `inicio` e `fim` geralmente usa o operador módulo.
\end{itemize}
\begin{center}
\begin{tikzpicture}
    \def\numsectors{6}
    \def\radius{1.7cm}
    \foreach \i [count=\j] in {A, B, C, D, ,} {
        \draw[thick, draw=lscText] ({360/\numsectors * (\j-1)}:\radius) arc ({360/\numsectors * (\j-1)}:{360/\numsectors * \j}:\radius) -- (0,0) -- cycle;
        \node[label_style] at ({360/\numsectors * (\j-0.5)}: \radius*0.8) {\huge \i};
        \node[label_style, font=\bfseries] at ({360/\numsectors * (\j-0.5)}: \radius*1.2) {\j};
    }
    \draw[pointer_style, thick, red] (34:\radius*1.5) -- node[above right,yshift=0.2cm] {início} (34:\radius*1.05);
    \draw[pointer_style, thick, red] (240:\radius*1.5) -- node[below left,yshift=-0.2cm] {fim} (240:\radius*1.05);
\end{tikzpicture}
\end{center}
\end{frame}


\begin{frame}{Comportamento}
\footnotesize
\begin{itemize}
  \item \textbf{Fila em vetor circular}
  \begin{itemize}
   \item Qual a condição para fila vazia? para fila cheia ?  Qual a condição inicial (quando a fila é criada)?
  \end{itemize}

 \end{itemize}

\begin{center}
\begin{tikzpicture}
    \def\numsectors{6}
    \def\radius{1.7cm}
    \foreach \i [count=\j] in {A, B, C, D, ,} {
        \draw[thick, draw=lscText] ({360/\numsectors * (\j-1)}:\radius) arc ({360/\numsectors * (\j-1)}:{360/\numsectors * \j}:\radius) -- (0,0) -- cycle;
        \node[label_style] at ({360/\numsectors * (\j-0.5)}: \radius*0.8) {\huge \i};
        \node[label_style, font=\bfseries] at ({360/\numsectors * (\j-0.5)}: \radius*1.2) {\j};
    }
    \draw[pointer_style, thick, red] (34:\radius*1.5) -- node[above right,yshift=0.2cm] {início} (34:\radius*1.05);
    \draw[pointer_style, thick, red] (240:\radius*1.5) -- node[below left,yshift=-0.2cm] {fim} (240:\radius*1.05);
\end{tikzpicture}
\end{center}

\end{frame}

\begin{frame}{Algoritmo}
\begin{center}
 \includegraphics[bb=0 0 644 536,scale=0.35]{figs/PseudoCodeFila1.png}
 % PseudoCodeFila1.png: 644x536 px, 72dpi, 22.72x18.91 cm, bb=0 0 644 536
\end{center}
\end{frame}


\begin{frame}{Algoritmo}
\begin{center}
 \includegraphics[bb=0 0 644 536,scale=0.35]{figs/PseudoCodeFila2.png}
 % PseudoCodeFila1.png: 644x536 px, 72dpi, 22.72x18.91 cm, bb=0 0 644 536
\end{center}
\end{frame}

\begin{frame}[fragile]{Implementação em C}

\centering\footnotesize
\begin{lstlisting}[language=C]
#include <stdio.h>
#include <stdlib.h>

#define MAX_SIZE 10

int queue[MAX_SIZE];
int inicio = 0;
int fim = -1;
int size = 0;//Tamanho

\end{lstlisting}
\end{frame}

\begin{frame}[fragile]{Implementação em C}

\centering\footnotesize
\begin{lstlisting}[language=C]
void enqueue(int value) {
    if (size == MAX_SIZE) {//Fila Cheia
        return;
    }
    fim = (fim + 1);
    queue[fim] = value;
    size++;
}
int dequeue() {
    if (inicio > fim) {//Fila Vazia
        return -1;
    }
    int value = queue[inicio];
    inicio = (inicio + 1);
    size--;
    return value;
}

\end{lstlisting}
\end{frame}


\begin{frame}[fragile]{Implementação em Python.}

\centering\footnotesize
\begin{lstlisting}[language=Python]
class Fila:
    def __init__(self, tamanho_maximo):
        self.tamanho_maximo = tamanho_maximo
        self.fila = []
        self.inicio = 0
        self.fim = -1

    def queue_empty(self):
        return self.fim < self.inicio

    def queue_isFull(self):
        return len(self.fila) == self.tamanho_maximo

    def print_front(self):
        if not self.queue_empty():
            print(self.fila[self.inicio])


\end{lstlisting}
\end{frame}


\begin{frame}[fragile]{Implementação em Python.}

\centering\footnotesize
\begin{lstlisting}[language=Python]
class Fila:

    def enqueue(self, valor):
        if not self.queue_isFull():
            self.fim += 1
            self.fila.insert(self.fim,valor)
    def dequeue(self):
        if not self.queue_empty():
            valor = self.fila[self.inicio]
            self.inicio += 1
            return valor


\end{lstlisting}
\end{frame}


\begin{frame}{Análise}
 \begin{block}{Análise da Complexidade}
 \begin{description}
  \item[\textit{Enqueue}] é uma operação básica, como a de atribuição, e atualização de índices. Logo é uma operação de tempo constante, $O(1)$ 
  \item[\textit{Dequeue}] $O(1)$
  \item [Consulta a Fila] Como esta operação percorre toda a estrutura de tamanho ``n'', seu tempo de execução é $O(n)$ 
 \end{description}

 \end{block}


\end{frame}

\begin{frame}{Aplicações}

\begin{itemize}
 \item Em computação gráfica, uma aplicação interessante é o algoritmo para colorir regiões gráficas, que consiste em varrer uma matriz de pontos, enfileirar os pontos conectados a esse e trocar a cor, após remover cada um da fila.
 \begin{center}
 \includegraphics[bb=0 0 165 46]{figs/AppFila.png}
 % AppFila.png: 686x191 px, 300dpi, 5.81x1.62 cm, bb=0 0 165 46
\end{center}

\item Cálculo de distância em Grafos.
\end{itemize}

\end{frame}

\section{Trabalho 1 - Estruturas Genéricas.}
\begin{frame}
\begin{itemize}
 \item Selecione uma questão e faça:

 \begin{enumerate}
  %\item Você pode implementar a operação de conjuntos dinâmicos \textit{INSERT} em uma \underline{lista simplesmente encadeada} em tempo \textbf{O(1)}? E a operação DELETE ?
  % \item Implemente uma pilha usando uma lista simplesmente encadeada As operações PUSH e POP ainda devem demorar o tempo $O(1)$
    \item Implemente uma pilha usando uma lista simplesmente encadeada.
    \item Implemente uma fila usando uma lista simplesmente encadeada.
%
%  \item Implemente uma fila usando uma lista simplesmente encadeada As operações ENQUEUE e DEQUEUE ainda devem demorar o tempo $O(1)$
 %\item Dê um procedimento não recursivo de tempo $\Theta(n)$ que inverta uma lista simplesmente encadeada de $n$ elementos. O procedimento só pode usar armazenamento constante além do necessário para a própria lista
%  \item Mostre como implementar uma fila usando duas pilhas. Analise o tempo de execução das operações nessa fila.
 \end{enumerate}
 \item Faça a lista usando Ponteiro ou autoreferência (Veja nos Livros).
 \end{itemize}


\end{frame}
\section{Referências Bibliográficas}
\begin{frame}[t,allowframebreaks]
\frametitle{Referências}
%\bibliographystyle{amsalpha}
%\bibliography{bibfile}
\begin{thebibliography}{*}

        \bibitem[A01]{A01} Lee K.D., Hubbard S. (2015) Trees. In: Data Structures and Algorithms with Python. Undergraduate Topics in Computer Science. Springer, Cham. Retrieved from \url{https://doi.org/10.1007/978-3-319-13072-9_6}
        
        \bibitem[A02]{A02}Hubbard, J. (2007). Schaum’s Outline sof Data Structures with Java. Retrieved from \url{http://www.amazon.com/Schaums-Outline-Data-Structures-Java/dp/0071476989}
        
        \bibitem[A03]{A03} Cormen, T. H., Leiserson, C. E., \& Stein, R. L. R. E. C. (2012). Algoritmos: teoria e prática. Retrieved from \url{https://books.google.com.br/books?id=6iA4LgEACAAJ}.
        
        \bibitem[A04]{A04} Ascenio, Ana Fernanda Gomes. Estrutura de dados: Algoritmos, análise da complexidade e implementações em Java e C++. São Paulo: Pearson Prentice Hall, 2010.
        
        \bibitem[A06]{A06} Szwarcfiter, Jayme Luiz. Estruturas de dados e seus algoritmos / Jayme Luiz Szwarcfiter, Lilian Markenzon. 3.ed. [Reimpr.]. - Rio de Janeiro : LTC, 2015.
        \bibitem[A07]{A07} Capítulo 20: Árvores — documentação Aprenda Computação com Python 3.0 2009.1. Retrieved from  \url{https://mange.ifrn.edu.br/python/aprenda-com-py3/capitulo_20.html}
\end{thebibliography}
\end{frame}
\section{}

%remover logo
\setbeamertemplate{logo}{}
\begin{frame}
\titlepage
\end{frame}

\end{document}
